\documentclass[../cours_1re.tex]{subfiles}
\begin{document}
\setcounter{chapter}{26}
\chapter{Complément : espaces vectoriels et projections}
\label{ch:27}
\begin{objectif}
Complément de transition vers la licence : généraliser les vecteurs, les produits scalaires et les projections.
Prérequis : chapitres 16 et 18 ; intégration du chapitre 25 pour le dernier exemple.
\end{objectif}
\section{Des vecteurs du plan aux listes de nombres}
Un vecteur de $\R^n$ est une liste ordonnée $(x_1,\ldots,x_n)$ de réels.
L'addition et la multiplication par un réel se font coordonnée par coordonnée.
Les règles de calcul se vérifient sur chaque coordonnée.
Plus généralement, un \emph{espace vectoriel réel} est un ensemble muni d'une
addition et d'une multiplication par un réel : l'addition est associative
et commutative, possède un neutre $0$ et un opposé pour chaque vecteur ;
la multiplication vérifie, pour tous réels $\lambda,\mu$ et vecteurs $u,v$,
\[
1u=u,\quad(\lambda\mu)u=\lambda(\mu u),\quad
\lambda(u+v)=\lambda u+\lambda v,\quad
(\lambda+\mu)u=\lambda u+\mu u.
\]
Les opérations restent dans l'ensemble considéré. Les polynômes réels
en fournissent un autre exemple ; le produit scalaire n'est pas requis
pour définir un espace vectoriel.

Une combinaison linéaire est une somme $\sum_{i=1}^r\lambda_i v_i$.
Une famille est libre si $\sum_i\lambda_iv_i=0$ force tous les coefficients
à être nuls ; elle est génératrice si tout vecteur est une telle combinaison.
Une base est une famille à la fois libre et génératrice.
Dans $\R^n$, les vecteurs $e_i$ ayant un $1$ à la position $i$ et des $0$
ailleurs forment une base : $x=\sum_ix_ie_i$, de manière unique.

\section{Applications linéaires et matrices}
Une application $T$ est linéaire si
$T(au+bv)=aT(u)+bT(v)$ pour tous $u,v,a,b$.
Elle est entièrement déterminée par les images des vecteurs d'une base :
$T(\sum_ix_ie_i)=\sum_ix_iT(e_i)$.
Une matrice range ces images en colonnes ; multiplier la matrice par
le vecteur des coordonnées calcule cette combinaison.
\begin{exemple}
$T(x,y)=(x+y,2x+2y)$ a pour matrice
$\begin{pmatrix}1&1\\2&2\end{pmatrix}$.
Son noyau, ensemble des vecteurs envoyés sur $0$, est
$\{(t,-t):t\in\R\}$.
Son image est la droite $\{(s,2s):s\in\R\}$.
Le noyau et l'image sont stables par combinaison linéaire, donc des
sous-espaces vectoriels. En particulier, $T$ ne permet pas de retrouver
un antécédent unique à partir de son image.
\end{exemple}

\section{Produit scalaire et distance}
Dans $\R^n$, on pose $\langle x,y\rangle=\sum_ix_iy_i$.
C'est une forme bilinéaire symétrique, positive, et
$\langle x,x\rangle=0$ seulement pour $x=0$.
La norme associée est $\|x\|=\sqrt{\langle x,x\rangle}$.
La preuve par le carré $\|x-ty\|^2\ge0$ du chapitre 16 démontre encore
Cauchy--Schwarz ; le développement de $\|x+y\|^2$ prouve l'inégalité triangulaire.
La distance $d(x,y)=\|x-y\|$ vérifie donc les propriétés attendues d'une distance.
Une famille est orthonormée si $\langle e_i,e_j\rangle$ vaut $1$ pour $i=j$
et $0$ sinon. Elle est libre : prendre le produit scalaire d'une combinaison
nulle avec chaque $e_j$ donne $\lambda_j=0$.

\section{Construire une base orthonormée}
Pour deux vecteurs libres $u,v$, poser
\[
e_1=\frac u{\|u\|},\qquad w=v-\langle v,e_1\rangle e_1,\qquad
e_2=\frac w{\|w\|}.
\]
On a $\langle w,e_1\rangle=0$. Le vecteur $w$ n'est pas nul, sinon
$v$ serait proportionnel à $u$. Les deux vecteurs $e_1,e_2$ sont donc
orthonormés et engendrent le même plan que $u,v$.
C'est la première étape du procédé de Gram--Schmidt.
\begin{exemple}
Pour $u=(1,1,0)$ et $v=(1,0,1)$,
$e_1=(1,1,0)/\sqrt2$ et
$w=(1/2,-1/2,1)$, donc $e_2=(1,-1,2)/\sqrt6$.
Le produit scalaire des numérateurs vaut $1-1+0=0$.
\end{exemple}

\section{Projection : la meilleure approximation}
Si $F$ possède une base orthonormée $(e_1,\ldots,e_r)$, poser
$p_F(x)=\sum_i\langle x,e_i\rangle e_i$.
Le vecteur $x-p_F(x)$ est orthogonal à chaque $e_j$, donc à tout $F$.
Pour tout $z\in F$, les deux termes de
$x-z=(x-p_F(x))+(p_F(x)-z)$ sont orthogonaux. Ainsi
\[
\|x-z\|^2=\|x-p_F(x)\|^2+\|p_F(x)-z\|^2.
\]
Le minimum est atteint uniquement pour $z=p_F(x)$.
Cette preuve établit à la fois l'optimalité et l'unicité, sans calcul différentiel.
En particulier, sur une droite $\R u$ avec $u\ne0$,
$p(x)=\langle x,u\rangle u/\|u\|^2$.

\section{Ajuster une droite par moindres carrés}
Pour approcher les points $(0,1),(1,2),(2,2)$ par $y=ax+b$, on minimise
\[
Q(a,b)=(b-1)^2+(a+b-2)^2+(2a+b-2)^2.
\]
Le vecteur des valeurs prédites est $a(0,1,2)+b(1,1,1)$.
À la projection, le résidu est orthogonal à ces deux vecteurs :
\[
5a+3b=6,\qquad 3a+3b=5.
\]
On obtient $a=1/2$, $b=7/6$. Le résidu « observation moins prédiction »
vaut $(-1/6,1/3,-1/6)$, bien orthogonal aux deux vecteurs.
L'identité de projection prouve que c'est l'unique minimum global, de valeur
$1/6$. Les équations dites normales ne sont pas seulement des recettes de calcul.

\section{Un espace de fonctions}
Sur les polynômes, $\langle P,Q\rangle=\int_0^1P(t)Q(t)\,dt$ est un produit
scalaire. La linéarité et la symétrie viennent de l'intégrale.
Si $P\ne0$, il n'est pas nul sur tout $[0,1]$ (un polynôme non nul a
un nombre fini de racines). Par continuité, $P^2$ est strictement positif
sur un petit intervalle ; son intégrale est donc positive.
On peut orthogonaliser $1,t,t^2,\ldots$ comme des vecteurs.
Par exemple $t-1/2$ est orthogonal à $1$ car son intégrale est nulle.

\begin{finchapitre}
\section{Utilité et prolongements}
Les coordonnées permettent de calculer ; les propriétés vectorielles
expliquent pourquoi le calcul fonctionne. Une projection minimise une erreur au carré.

En licence, matrices, valeurs propres et espaces euclidiens systématiseront ces
idées. En master, les espaces de Hilbert demanderont en plus une propriété de
complétude ; elle ne découle pas de la seule présence d'un produit scalaire.
\begin{retenir}
Une matrice représente une application linéaire dans des bases.
L'orthogonalité du résidu caractérise une projection et prouve l'optimalité des moindres carrés.
\end{retenir}
\end{finchapitre}
\end{document}
